\subsection{Non-Committing Encryption}
\label{app:nce}

We give the key-first experiment of Hemenway et al.\
\cite[Definition~2]{HemenwayEtAl2016QuasiOptimalNCE}
in the notation of Section~\ref{sec:nce}. For random-tape exposure see also
Brakerski et al.~\cite[full version, Definition~11]{BrakerskiEtAl2020ConstantNCE}
and Damg{\aa}rd and Nielsen~\cite[full version, Section~4]{DamgardNielsen2000NCE}.
Correctness requires, for every efficient algorithm $\mathcal B$ that
chooses a message $m\in\bits^\ell$ after seeing $\pk$,
\[
  \Pr\bigl[\Dec(\sk,c)=m\bigr]\geq 1-\negl(\lambda),
\]
where $(\pk,\sk)\sample\KeyGen(1^\lambda)$,
$m\sample\mathcal B(1^\lambda,\ell,\pk)$, and
$c\sample\Enc(\pk,m)$.  The probability includes the decryption coins.  For every fixed polynomial
bound on the public message length, the correctness error is uniform
over all lengths within that bound.

The receiver and sender tapes $\rho_R$ and $\rho_S$ are uniform strings
of their prescribed polynomial lengths. A receiver that decrypts also
retains its independent decryption coins. The experiment has no private decryption step before corruption.

The experiments $\NCEReal_{\mathcal A}(1^\lambda)$ and
$\NCEIdeal_{\mathcal A,\mathcal S}(1^\lambda)$ give $1^\lambda$ and
$\ell$ to the efficient adversary $\mathcal A$ and offer these commands:
\begin{description}
  \item[$\GetPublicKey()$.]
    In the real experiment, sample $\rho_R$ uniformly and compute
    \[
      (\pk,\sk)\gets\KeyGen(1^\lambda;\rho_R).
    \]
    Retain $(\sk,\rho_R)$ and return $\pk$.
    In the ideal experiment, compute
    \[
      (\widetilde{\pk},\widetilde c,\SimState)\sample\SimNCE(1^\lambda),
    \]
    retain $(\widetilde c,\SimState)$, and return $\widetilde{\pk}$.
    The precomputed ciphertext remains private until $\EncryptMessage$.
  \item[$\EncryptMessage(m)$.]
    Store $m\in\bits^\ell$, which may depend on the public key.
    In the real experiment, sample $\rho_S$ uniformly and return
    $c\gets\Enc(\pk,m;\rho_S)$.
    In the ideal experiment, return the stored ciphertext $\widetilde c$.
    The simulator is not given $m$ at this step.
  \item[$\CorruptSender()$.]
    In the real experiment, return $(m,\rho_S)$.
    In the ideal experiment, return $(m,\widetilde\rho_S)$, using the
    cached explanation described below.
  \item[$\CorruptReceiver()$.]
    In the real experiment, return $(\sk,\rho_R)$.
    In the ideal experiment, compute
    $(\pk',\widetilde{\sk})\gets
      \KeyGen(1^\lambda;\widetilde\rho_R)$
    from the cached explanation, and return
    $(\widetilde{\sk},\widetilde\rho_R)$.
  \item[$\Finish(b)$.]
    Terminate and output $b\in\bits$.
\end{description}
On the \emph{first} corruption command in the ideal experiment, give
$m$ to the simulator, compute
\[
  (\widetilde\rho_R,\widetilde\rho_S)
    \sample\ExplainNCE(\SimState,m),
\]
and cache this pair for both corruption commands.

The commands $\GetPublicKey$ and $\EncryptMessage$ are accepted once
each, in that order.  Each corruption command is accepted at most once,
only after $\EncryptMessage$, and the two corruptions may occur in either
order.  The adversary may omit either or both.  Malformed, out-of-order,
or repeated commands return $\bot$ without changing state.  $\Finish$
may be called at any time; termination without an output bit means
$\Finish(0)$.  Allowing the adversary to request the two parts of one
explanation at different times is just post-processing of the joint
exposure in the cited definitions.

\begin{definition}[Non-committing security]
\label{def:nce}
The scheme has non-committing security with tape exposure if there exists
an efficient simulator $\mathcal S=(\SimNCE,\ExplainNCE)$ such that,
for every efficient adversary $\mathcal A$ and polynomial $b$, there is a
negligible function $\varepsilon_{\mathcal A,b}$ satisfying
\[
  \left|
    \Pr[\NCEReal_{\mathcal A}(1^\lambda)=1]
    -\Pr[\NCEIdeal_{\mathcal A,\mathcal S}(1^\lambda)=1]
  \right|\leq\varepsilon_{\mathcal A,b}(\lambda).
\]
The bound holds uniformly for every public length
$0\leq\ell\leq b(\lambda)$.  The same algorithms take $\ell$ as a public
parameter; the game fixes it before invoking $\mathcal A$.
\end{definition}

\begin{definition}[Sender and receiver restrictions]
\label{def:sender-nce}
\label{def:receiver-nce}
\emph{Sender non-committing encryption} (SNCE) requires the same
indistinguishability only for adversaries that never call
$\CorruptReceiver$.  \emph{Receiver non-committing encryption} (RNCE)
requires it only for adversaries that never call $\CorruptSender$.
Each restriction permits its own simulator.  Full NCE implies both by
discarding the unused part of the explanation.  The receiver-only
notion originates with Jarecki and Lysyanskaya~
\cite{JareckiLysyanskaya2000AdaptiveThreshold}.
\end{definition}

In the full game the explained tapes must, except with negligible
probability, satisfy
\[
  \KeyGen(1^\lambda;\widetilde\rho_R)
    =(\widetilde{\pk},\widetilde{\sk}),\qquad
  \Enc(\widetilde{\pk},m;\widetilde\rho_S)=\widetilde c.
\]
They must also have the correct joint computational distribution, not
merely satisfy these equations.  Decrypting $\widetilde c$ with the
explained secret key and fresh coins recovers $m$ except with negligible
probability.  Explained tapes describe honest computations; they need
not be the tapes actually used internally by the simulator.
